\subsection{非齐次线性方程的积分}

求非齐次线性方程

\[
{\frac {d \mathbf {x}}{d t}} = A (t). \mathbf {x} + \mathbf {b} (t)\tag{2}
\]

的积分可化归为求相伴的齐次方程

\[
{\frac {d \mathbf {x}}{d t}} = A (t). \mathbf {x}\tag{4}
\]

的积分并计算一个原函数。采用 IV 第 179 页定理 2 的记号，令 \(\mathbf{x} = C(t, t_0).\mathbf{z}\) ，于是由 IV 第 179 页的第二公式 (6)，\(\mathbf{z} = C(t_0, t).\mathbf{x}\) ；若 \(\mathbf{x}\) 是 (2) 的积分，则 \(\mathbf{z}\) 是方程 \(\frac{d}{dt}\big(C(t, t_0).\mathbf{z}\big) = A(t)C(t, t_0).\mathbf{z} + \mathbf{b}(t)\) 的积分；由于双线性映射

\[
(U, \mathbf {y}) \mapsto U. \mathbf {y}
\]

（从 \(\mathcal{L}(E) \times E\) 到 E 中）是连续的（《一般拓扑学》，X，第 297 页，命题 6），z 有导数（除去 J 的一个可数子集外），并且由双线性函数的求导公式（I，第 6 页，命题 3）有

\[
\frac {d}{d t} \big (C (t, t _ {0}). \mathbf {z} \big) = \frac {d C (t , t _ {0})}{d t}. \mathbf {z} + C (t, t _ {0}). \frac {d \mathbf {z}}{d t} = A (t) C (t, t _ {0}). \mathbf {z} + C (t, t _ {0}). \frac {d \mathbf {z}}{d t}
\]

（把 \(dC(t, t_{0})/dt\) 换成 \(A(t)C(t, t_{0})\) ，依据 (5)（IV，第 179 页））。于是 z 的方程化归为 \(C(t, t_{0}).d\mathbf{z}/dt = \mathbf{b}(t)\) ，或者再化归为

\[
\frac {d \mathbf {z}}{d t} = C (t _ {0}, t). \mathbf {b} (t)\tag{8}
\]

这是由 IV 第 179 页的第二公式 (6) 得到的。现在，方程 (8) 的右端是 J 上的规则函数，它是由把规则函数 U 与 y 代入连续双线性函数 U.y 而得到的（参看 II，第 55 页，推论 2）；于是方程 (8) 有唯一一个取值 \(x_{0}\) 于 \(t_{0}\) 的积分，它由下列公式给出

\[
\mathbf {z} (t) = \mathbf {x} _ {0} + \int_ {t _ {0}} ^ {t} C (t _ {0}, s). \mathbf {b} (s) d s.\tag{9}
\]

由于有 \(C(t, t_{0})\) · \(\int_{t_{0}}^{t} C(t_{0}, s)\) · \(\mathbf{b}(s) \, ds = \int_{t_{0}}^{t} C(t, t_{0}) C(t_{0}, s)\) · \(\mathbf{b}(s) \, ds\) （II，第 59 页，公式 (9)），便得到（考虑到 IV 第 179 页的第一公式 (6)）下列结果：

命题 3. 采用定理 2（IV，第 179 页）的记号，对 \((t_0, \mathbf{x}_0)\) 的每个点 \(\mathbf{J} \times \mathbf{E}\)，线性方程 (2) 的定义在 \(\mathbf{J}\) 上且等于 \(\mathbf{x}_0\) 于 \(t_0\) 的积分由下列公式给出

\[
\mathbf {u} (t) = C (t, t _ {0}). \mathbf {x} _ {0} + \int_ {t _ {0}} ^ {t} C (t, s). \mathbf {b} (s) d s.\tag{10}
\]

导致公式 (10) 的方法，即取函数 z 作为新的未知函数的方法，常称为“常数变易法”。
% semantic-fragments: legacy-input-seam
\relax{}
